% !PRINTED 92
\weberpage{92}
Eine dritte Art der Entwickelung, in der die Symmetrie der
Funktionen in bezug auf die beiden Variablen $v$, $a$ zum Ausdruck
kommt, ergibt sich, wenn man die Brüche, die in den vorigen
Entwickelungen auftraten, nach Potenzen von $q$ entwickelt. So
erhält man
\[
\begin{gathered}
\frac{q^m}{e^{-2\pi i v}-q^m} = \Sigma\, q^{m\nu} e^{2\pi i \nu v}, \\
\frac{q^m}{e^{+2\pi i v}-q^m} = \Sigma\, q^{m\nu} e^{-2\pi i \nu v}, \\
\nu = 1,\, 2,\, 3,\, \ldots
\end{gathered}
\]
und diese Entwickelungen gelten, solange der absolute Wert
von $q\, e^{2\pi i v}$ ein echter Bruch ist, also, wenn $v = v' + i v''$ gesetzt
wird, solange
\[
-\omega'' < v < \omega'' % [sic] Printed v; the imaginary-part bound mathematically requires v''.
\]
ist. Hiernach ergibt sich aus (14)
\[
\begin{gathered}
\frac{\vartheta'_{11}\,\vartheta_{11}(v+a)}{\pi\,\vartheta_{11}(v)\,\vartheta_{11}(a)} = \mathrm{cotg}\,\pi v + \mathrm{cotg}\,\pi a \\
{}+ 4\Sigma\, q^{\frac{mm'}{2}} \sin(m a + m' v)\pi,
\end{gathered}
\]
worin $m$, $m'$ voneinander unabhängig die \gesperrt{geraden} Zahlen
2, 4, 6, 8, \ldots\ durchlaufen. Diese Formel ist gültig für \gesperrt{reelle}
$a$ und $v$ und gilt darüber hinaus noch, solange der imaginäre
Teil, sowohl von $v$ als von $a$, absolut kleiner ist als $\omega''$.

In der Tabelle III sind die 16 Formeln dieser Art zusammengestellt\textsuperscript{1}).

\textsuperscript{1}) Solche Entwickelungen sind von \gesperrt{Jacobi} (sur la rotation d'un corps)
(ges. Werke Bd.\ II), hierauf von \gesperrt{Hermite} (Annales de l'école normale 1885)
und von \gesperrt{Kronecker} (Berliner Akademie 1885) betrachtet. Vgl.\ auch die
Straßburger Dissertation von L.\ \gesperrt{Vockerodt} (1905).
